\documentclass[11pt]{article}

\usepackage[margin=1in]{geometry}
\usepackage{amsmath,amssymb,amsthm}
\usepackage{graphicx}
\usepackage[hidelinks]{hyperref}

\newtheorem{theorem}{Theorem}
\newtheorem{lemma}{Lemma}
\newtheorem{remark}{Remark}

\newcommand{\Lip}{\operatorname{Lip}}
\newcommand{\Free}{\mathcal F}
\newcommand{\Ball}{B}
\newcommand{\norm}[1]{\left\lVert #1 \right\rVert}

\title{Metric Wedge Obstructions in a Lipschitz-Free Uniform Property \texorpdfstring{$(S)$}{(S)} Question}
\author{Automated Math Discovery Smoke Test}
\date{June 12, 2026}

\begin{document}
\maketitle

\section*{Source and Open Question}

This packet concerns the Lipschitz-free item in the open-problem section
of:

\begin{quote}
William B. Johnson and Tomasz Kania, \emph{Uniform Property (S)},
arXiv:2602.09106.
\end{quote}

The source location is page 19 of the paper.

\begin{center}
\includegraphics[width=\textwidth]{figures/open_problem_crop.png}
\end{center}

The paper defines property $(S)$ as follows: a normed space $X$ has
property $(S)$ if for every distinct $x,y\in S_X$ and every $a>0$ there is
$z\in X$ with $\norm{z}<a$ such that
\[
  \norm{x+z}\neq \norm{y+z}.
\]
Uniform property $(S)$ is the corresponding quantitative strengthening
using the modulus
\[
  U_X(d;a)=
  \inf_{\substack{x,y\in S_X\\ \norm{x-y}\ge d}}
  \sup_{\norm{z}\le a}
  \left|\norm{x+z}-\norm{y+z}\right|.
\]
Uniform property $(S)$ implies property $(S)$.

\section*{Metric Wedges}

Let $(A,0)$ and $(B,0)$ be pointed metric spaces. Their metric wedge
$M=A\vee B$ is the pointed metric space obtained by identifying the two
base points and declaring that
\[
  d(a,b)=d_A(a,0)+d_B(0,b)
  \qquad(a\in A,\ b\in B).
\]
The original metrics are kept on each piece.

\begin{lemma}
For a metric wedge $M=A\vee B$ there is a canonical isometric
identification
\[
  \Free(M) \cong \Free(A)\oplus_1\Free(B).
\]
\end{lemma}

\begin{proof}
Restriction to the two pieces gives an isometric identification
\[
  \Lip_0(M)=\Lip_0(A)\oplus_\infty \Lip_0(B).
\]
Indeed, if $f_A\in\Lip_0(A)$ and $f_B\in\Lip_0(B)$, the function obtained
by gluing $f_A$ and $f_B$ at the base point has Lipschitz constant
\[
  \max\{\Lip(f_A),\Lip(f_B)\}.
\]
The only cross-piece estimate needed is
\[
  |f_A(a)-f_B(b)|
  \le \Lip(f_A)d_A(a,0)+\Lip(f_B)d_B(b,0)
  \le \max\{\Lip(f_A),\Lip(f_B)\}d(a,b).
\]
Taking the canonical predual generated by point evaluations gives
$\Free(M) \cong \Free(A)\oplus_1\Free(B)$.
\end{proof}

\section*{Candidate Partial Result}

\begin{theorem}
Let $M=A\vee B$ be a metric wedge of pointed metric spaces. If $A$ is
finite and $|A|\ge 3$, then $\Free(M)$ fails property $(S)$, and
therefore fails uniform property $(S)$.
\end{theorem}

\begin{proof}
By the lemma,
\[
  \Free(M)=\Free(A)\oplus_1\Free(B).
\]
Write $Y=\Free(A)$ and $Z=\Free(B)$.

Since $A$ is finite, the unit ball of $Y$ is the centrally symmetric
polytope
\[
  \Ball_Y=
  \operatorname{conv}
  \left\{
  \frac{\delta_u-\delta_v}{d(u,v)}:
  u,v\in A,\ u\neq v
  \right\}.
\]
Also $\dim Y=|A|-1\ge 2$. Hence $\Ball_Y$ has a positive-dimensional
facet. Choose two distinct points $x,y$ in the relative interior of such
a facet, and let $\varphi$ be a norm-one functional exposing that facet.

Because $\Ball_Y$ is a polytope, there are finitely many norm-one
functionals $\varphi_1,\ldots,\varphi_N$ such that
\[
  \norm{u}_Y=\max_{1\le j\le N}\varphi_j(u)
  \qquad(u\in Y),
\]
with one of the $\varphi_j$ equal to $\varphi$. Since $x$ and $y$ lie in
the relative interior of the facet exposed by $\varphi$, all other facet
functionals are strictly smaller than $\varphi$ at both $x$ and $y$.
By finiteness and continuity there is $\eta>0$ such that whenever
$\norm{u}_Y<\eta$,
\[
  \norm{x+u}_Y=\varphi(x+u)
  \quad\hbox{and}\quad
  \norm{y+u}_Y=\varphi(y+u).
\]
Consequently,
\[
  \norm{x+u}_Y-\norm{y+u}_Y
  =
  \varphi(x+u)-\varphi(y+u)
  =
  \varphi(x)-\varphi(y)
  =
  0
\]
for every $\norm{u}_Y<\eta$.

Now view $x$ and $y$ as elements of $Y\oplus_1 Z$. They remain distinct
unit vectors. If $w=(u,v)\in Y\oplus_1 Z$ and $\norm{w}<\eta$, then
$\norm{u}_Y<\eta$, and therefore
\[
  \norm{x+w}
  =
  \norm{x+u}_Y+\norm{v}_Z
  =
  \norm{y+u}_Y+\norm{v}_Z
  =
  \norm{y+w}.
\]
Thus the pair $x,y\in S_{\Free(M)}$ cannot be separated by any
perturbation of norm $<\eta$. This violates property $(S)$, and hence
also uniform property $(S)$.
\end{proof}

\section*{Verification Notes}

No computational verifier was used for this packet. The proof is a
structural Banach-space argument. The main review points are:

\begin{itemize}
\item the standard metric-wedge identity
  $\Free(A\vee B)=\Free(A)\oplus_1\Free(B)$;
\item the polyhedral finite-free-space obstruction from the finite metric
  subcase;
\item the preservation of the obstruction in an $\ell_1$-sum, because the
  second coordinate contributes the same term to both perturbed norms.
\end{itemize}

\section*{Limitations and Novelty Check}

This packet is only a partial obstruction for the Lipschitz-free open
problem. It does not characterize all pointed metric spaces $M$ for which
$\Free(M)$ has uniform property $(S)$.

The result strengthens the earlier finite-metric packet only when the
second wedge component is nontrivial, especially when it is infinite.
The novelty risk is moderate: both ingredients, the metric-wedge
decomposition of Lipschitz-free spaces and the finite polyhedral
obstruction, are standard. The packet records the combined obstruction as
a convenient explicitly checked subcase of the source paper's
characterization problem.

The same-source check found no nearby answer to this wedge obstruction in
the source paper. The local run registry contains only the finite-metric
free-space subcase, not this wedge extension.

\end{document}
